%% ma: L11-B16-0009
%% lop: 11
%% chuong: 5
%% bai: 16
%% dang: 11.16.5
%% loai: TL
%% muc: VD
%% dap_an: "-"
%% nguon: "tai-lieu-hoc-tap-mon-toan-lop-11-theo-chuong-trinh-giao-duc-pho-thong-2018.pdf (Lớp Toán Thầy Duy, Phạm Lê Duy, THPT Chu Văn An), Chương 3 Bài 2, Dạng 3, Ví dụ 9 (tr.357)"
%% kien_thuc: "Lớp 11 Bài 16 (giới hạn của hàm số)"
%% trang_thai: cho-duyet
%% ly_do: "Dạng mới đề xuất (11.16.x), chờ giáo viên duyệt dạng; tài liệu gốc không có lời giải mẫu, đáp án do kho tự giải và kiểm bằng sympy."
%% ghi_chu: "Ý d) nhìn như ghi trong đề gốc (mẫu 2 - x^2 - 3x không triệt tiêu tại 1 nên đáp số 0); nghi đề gốc muốn mẫu triệt tiêu, nhờ giáo viên đối chiếu."
\begin{ex}
Tính các giới hạn sau.

a) $\lim\limits_{x\to-4^+}\dfrac{x^3-16x}{|x+4|}$;

b) $\lim\limits_{x\to-4^-}\dfrac{\sqrt{x^2-16}}{|x+4|}$;

c) $\lim\limits_{x\to2^-}\dfrac{|2-x|}{2x^2-5x+2}$;

d) $\lim\limits_{x\to1^-}\dfrac{|x-1|}{2-x^2-3x}$.
\loigiai{\sloppy a) Khi $x\to-4^+$: $|x+4|=x+4$, $x^3-16x=x(x-4)(x+4)$, nên biểu thức $=x(x-4)\to32$.

b) Khi $x\to-4^-$: $x+4<0$, $|x+4|=-(x+4)$. Biểu thức $=\dfrac{\sqrt{(4-x)(-x-4)}}{-(x+4)}=\sqrt{\dfrac{4-x}{-(x+4)}}$. Tử $4-x\to8>0$ và $-(x+4)\to0^+$, nên giới hạn bằng $+\infty$.

c) Khi $x\to2^-$: $|2-x|=2-x$ và $2x^2-5x+2=(2x-1)(x-2)$, nên biểu thức $=\dfrac{-1}{2x-1}\to-\dfrac13$.

d) Mẫu $2-x^2-3x\to-2\ne0$ và tử $|x-1|\to0$, nên giới hạn bằng $0$.}
\end{ex}
